% p005_c 的电脑重画（② 遏止电压 U_c 测定电路：左板(阳极,+)、右板(阴极,-)，电流表 A 示数 I=0，电压表 V，滑动变阻器与电池）
\begin{tikzpicture}[line width=0.7pt, >={Stealth[length=2.2mm]}]
  % 左侧导线与极板
  \draw (0.8,3.06) -- (0,3.06) -- (0,0) -- (1.83,0);
  \draw (0.65,3.39) -- (1.02,2.52);
  \node at (0.45,3.41) {$+$};
  \node at (0.47,3.03) {$+$};
  \node at (0.61,2.66) {$+$};
  % 右极板、负号、导线与电流表
  \draw (2.78,3.44) -- (2.78,2.39);
  \node at (3.0,3.3) {$-$};
  \node at (2.97,2.79) {$-$};
  \node at (2.95,2.51) {$-$};
  \draw (2.78,3.07) -- (3.35,3.07) -- (3.35,0.97) -- (1.65,0.97);
  \draw (3.35,2.03) circle (0.29);
  \node at (3.35,2.03) {A};
  \node at (4.15,2.0) {$I=0$};
  % 电压表（导线穿过圆圈）
  \draw (0,1.27) -- (3.35,1.27);
  \draw (1.71,1.3) circle (0.34);
  \node at (1.71,1.3) {V};
  % 滑动变阻器
  \draw (0,0.55) -- (1.1,0.55);
  \draw (1.1,0.41) rectangle (1.99,0.7);
  \draw (1.99,0.55) -- (3.35,0.55) -- (3.35,0) -- (2.02,0);
  \draw[->] (1.55,0.95) -- (1.55,0.7);
  % 电池
  \draw (1.83,-0.27) -- (1.83,0.27);
  \draw (2.02,-0.16) -- (2.02,0.16);
  % 电场、电子、光
  \draw[->, very thick] (1.67,3.5) -- (2.3,3.5);
  \node at (1.9,3.72) {$E$};
  \node at (1.2,3.17) {$e^-$};
  \draw[->] (1.63,2.95) -- (2.18,2.97);
  \draw[->] (3.05,3.95) -- (2.55,3.57);
  \node at (3.35,3.85) {$h\nu$};
  % 手写标注
  \node at (-0.6,3.23) {阳极};
  \node at (4.4,3.23) {阴极};
\end{tikzpicture}
